\documentclass[a4paper]{article}

\usepackage{a4wide}
\usepackage[utf8]{inputenc}
\usepackage[ngerman]{babel}
\usepackage{fancyhdr}
\usepackage{graphicx}
\usepackage[parfill]{parskip}
\pagestyle{fancy}
\topmargin -5mm
\headheight 20mm

% include logo in document head
\chead{\includegraphics[width=50mm]{logo.pdf}}

% bottom line
\renewcommand{\footrulewidth}{0.4pt}

% change standard font
\usepackage[T1]{fontenc}
\newcommand{\changefont}[3]{
\fontfamily{#1} \fontseries{#2} \fontshape{#3} \selectfont}

\newcommand{\TOP}[1]{\item \textbf{#1}\par}

\begin{document}
\changefont{cmss}{m}{n} %computer modern sans serif

\begin{center}
\textbf{\Large Protokoll zum Gruppentreffen vom 12. April 2014}
\end{center}

\textbf{Beginn:} 15:00 Uhr, \textbf{Ende:} 16:45 Uhr\\
\textbf{Anwesende:} Alex, Dennis, Philipp, Robert, Kilian

\vspace{5mm}

\begin{enumerate}

\TOP{Namensfindung}
Die Namensentscheidung fiel letzendlich auf \textit{AlphaLabs}. Alpha da wir Projektgruppe 1 sind.

\TOP{Logodesign}
Ein passendes Logo wird in den nächsten Tagen noch von Alex entworfen und dann in die Dokumentvorlage eingebunden.

\TOP{Framework}
Als Basistechnologie legen wir uns HTML5, CSS und Javascript bereit. Auf genauere Details und passende Frameworks (beispielsweise Pixi.js\footnote{http://www.pixijs.com/} oder Phaser\footnote{http://phaser.io/}) fällt die Entscheidung sobald die Konzeptidee ausgewählt ist.

\TOP{Konzeptideen}
Wir haben uns auf zwei mögliche Konzepte für das Spiel geeinigt. Ein Lernspiel zum Thema Mülltrennung und ein Geographieratespiel. Die Ausführungen hierzu sind separat.

\end{enumerate}



\end{document}
